\documentclass[11pt]{article}
\usepackage[a4paper,margin=25mm]{geometry}
\usepackage{amsmath,amssymb,amsthm,booktabs,array}
\usepackage[T1]{fontenc}
\usepackage{lmodern}
\usepackage[colorlinks=true,linkcolor=blue,urlcolor=blue,citecolor=blue]{hyperref}
\newtheorem{theorem}{Theorem}
\newtheorem{lemma}[theorem]{Lemma}
\newtheorem{corollary}[theorem]{Corollary}
\DeclareMathOperator{\ex}{exp}
\title{Proof of Hanna's square-exponent congruence\texorpdfstring{\\}{ }and an extension of Fried's divisor-sum formula}
\author{Haobo Ma\\\small ChronoAI Pte Ltd, Singapore
\and Wenlin Zhang\\\small National University of Singapore\\\small ChronoAI Pte Ltd, Singapore}
\date{18 November 2026}
\begin{document}
\maketitle
\begin{abstract}
For Paul D. Hanna's conjecture in OEIS A226838, revision 14, FORMULA,
we prove that every factorial-normalized coefficient of
$\exp(\sum_{m\geq1}x^{m^2}/m^m)$ is an integer congruent to one modulo six
by Carlitz's coefficient recurrence, and extend Sela Fried's
Theorem~3 (2025, page~3) on Yanev's divisor-sum identity to all positive exponents.
The same argument proves the modulo-three and modulo-thirty conjectures
in A226839 and A226890, with all three uniform moduli optimal.
A necessary and sufficient coefficient criterion gives a sharp weighted
extension; the linked divisor-sum entry's elementary identities also follow.
\end{abstract}

\section{The question and an integral recurrence}

Hanna asks why the sequence defined by
\begin{equation}\label{eq:square}
 A(x)=\ex\!\left(\sum_{m\geq1}\frac{x^{m^2}}{m^m}\right)
      =\sum_{n\geq0}a_n\frac{x^n}{n!}
\end{equation}
satisfies $a_n\equiv1\pmod6$ for every $n\geq0$~\cite{square}.
Integrality is part of the assertion: the defining exponential is initially
a formal power series over $\mathbb Q$.
The source records a numerical check through $n=1024$.
We give an all-index proof and then treat its two cross-references~\cite{triangle,weighted}.
We use Carlitz's Bell-polynomial recurrence~\cite[\S2, (2.1)--(2.2)]{carlitz};
its useful input here is an elementary factorial divisibility.
For the divisor-sum companion we use prime-power geometric sums, extending
Fried's proof for exponent two~\cite[Theorem~3, p.~3]{fried} to every positive
exponent in Yanev's conjecture~\cite[FORMULA, 24 August 2017]{divisors}.

For a series $F(x)$, write $[x^n]F(x)$ for its coefficient of $x^n$,
and set $0!=1$. All series below are formal series in $\mathbb Q[[x]]$;
their exponents have zero constant term, so their formal exponentials
are defined coefficientwise by finite sums.
The following recurrence avoids reducing rational factorial denominators
modulo a composite integer.

\begin{lemma}\label{lem:transfer}
Let $H(x)=\sum_{j\geq1}b_jx^j/j!$ with $b_j\in\mathbb Z$, and write
$\ex(H(x))=\sum_{n\geq0}u_nx^n/n!$.
Then all $u_n$ are integers. For $M\geq1$ and $b_1=t\in\mathbb Z$,
\[
 \bigl(\forall n\geq0,\ u_n\equiv t^n\pmod M\bigr)
 \quad\Longleftrightarrow\quad
 \bigl(\forall j\geq2,\ M\mid b_j\bigr),
\]
with $t^0=1$ even when $t=0$.
\end{lemma}
\begin{proof}
Put $U=\ex(H)$. Formal differentiation gives $U'=H'U$.
Comparing the coefficient of $x^n$ and multiplying by $n!$ yields
\begin{equation}\label{eq:recurrence}
 u_0=1,\qquad
 u_{n+1}=\sum_{j=0}^n\binom{n}{j}b_{j+1}u_{n-j}
 \quad(n\geq0).
\end{equation}
These are the coefficients of the given rational exponential, rather than
an independently postulated integer sequence. Induction in
\eqref{eq:recurrence} proves integrality. Modulo $M$, every summand
with $j\geq1$ vanishes, leaving $u_{n+1}\equiv t u_n$.
Induction from $u_0=1$ proves the forward construction.
Conversely, suppose $u_n\equiv t^n$ for all $n$. For $n\geq1$ isolate
the last summand $b_{n+1}u_0=b_{n+1}$ in \eqref{eq:recurrence}.
If $M\mid b_2,\ldots,b_n$, the remaining terms give
$b_{n+1}\equiv u_{n+1}-t u_n\equiv0\pmod M$.
Induction, starting with $n=1$, proves necessity.
\end{proof}

\section{Square exponents}

\begin{lemma}\label{lem:six}
For every integer $m\geq2$,
\begin{equation}\label{eq:divsix}
 6m^m\mid(m^2)!.
\end{equation}
\end{lemma}
\begin{proof}
For $m=2,3$, respectively, the quotients $(m^2)!/m^m$ are $6$ and
$13440$, both divisible by six.
For $m\geq4$, $(m-1)!$ is divisible by six, so $m!=6m c$ for an integer $c$.
The number of ordered decompositions of a set of $m^2$ elements into
$m$ disjoint subsets of size $m$ is the integer
\[
 \frac{(m^2)!}{(m!)^m}
   =\binom{m^2}{m}\binom{m^2-m}{m}\cdots\binom{m}{m}.
\]
Consequently $(m!)^m\mid(m^2)!$.
Since $(m!)^m=6^m m^m c^m$ is divisible by $6m^m$, \eqref{eq:divsix} follows.
\end{proof}

\begin{theorem}\label{thm:square}
The coefficients in \eqref{eq:square} satisfy $a_n\in\mathbb Z$ and
$a_n\equiv1\pmod6$ for every integer $n\geq0$.
\end{theorem}
\begin{proof}
Write the exponent in \eqref{eq:square} as $\sum_{j\geq1}b_jx^j/j!$.
There is at most one positive $m$ with $j=m^2$, and
\[
 b_j=\begin{cases}(m^2)!/m^m,&j=m^2,\ m\geq1,\\
                   0,&j\text{ is not a positive square}.
       \end{cases}
\]
Thus $b_1=1$, and Lemma~\ref{lem:six} shows that every $b_j$ with $j\geq2$
is an integer divisible by six. Lemma~\ref{lem:transfer}, with $M=6$ and
$t=1$, proves the assertion, including $a_0=1$.
\end{proof}

The modulus six is the largest uniform modulus for this sequence:
$a_4=7$, as is seen from the terms $\ex(x)(1+x^4/4+O(x^8))$,
so any integer dividing every $a_n-1$ must divide $a_4-1=6$.
\section{The two cross-referenced conjectures}

Define the triangular numbers $T_m=m(m+1)/2$ and the two sequences
\begin{align}
 B(x)&=\ex\!\left(\sum_{m\geq1}\frac{x^{T_m}}{m!}\right)
      =\sum_{n\geq0}v_n\frac{x^n}{n!},\label{eq:triangular}\\
 C(x)&=\ex\!\left(\sum_{m\geq1}\sigma_m(m)\frac{x^{m^2}}{m^m}\right)
      =\sum_{n\geq0}z_n\frac{x^n}{n!},\label{eq:weighted}
\end{align}
where $\sigma_m(m)=\sum_{d\mid m,\ d>0}d^m$ is the sum of the $m$th powers
of the positive divisors of $m$.
These are exactly the definitions in A226839 and A226890.

\begin{corollary}\label{cor:triangular}
For every $n\geq0$, $v_n\in\mathbb Z$ and $v_n\equiv1\pmod3$.
\end{corollary}
\begin{proof}
The factorial-normalized inner coefficient at degree $T_m$ is
$T_m!/m!$, an integer, because $T_m\geq m$.
It equals one for $m=1$ and three for $m=2$.
For $m\geq3$, the product $(m+1)(m+2)\cdots T_m$ contains
at least three consecutive integers, since $T_m-m=m(m-1)/2\geq3$.
It is therefore divisible by three. All other inner coefficients are zero.
Lemma~\ref{lem:transfer} applies with $M=3$ and $t=1$.
\end{proof}

\begin{corollary}\label{cor:weighted}
For every $n\geq0$, $z_n\in\mathbb Z$ and $z_n\equiv1\pmod{30}$.
\end{corollary}
\begin{proof}
Write $q_m=(m^2)!/m^m$. We first show $30\mid q_m$ for $m\geq3$.
For $m\geq6$, $30\mid(m-1)!$, and hence $30m\mid m!$.
The ordered-decomposition argument in Lemma~\ref{lem:six} then gives
$30m^m\mid(m^2)!$.
For the remaining cases $m=3,4,5$, the following exact valuations suffice.
Here $\nu_p(r)$ is the largest exponent $e$ such that $p^e\mid r$
for a positive integer $r$; factorial valuations count the multiples of
$p,p^2,\ldots$ in the factorial product.
\[
\begin{array}{c|ccc}
 m&\nu_2(q_m)&\nu_3(q_m)&\nu_5(q_m)\\\hline
 3&7&1&1\\
 4&7&6&3\\
 5&22&10&1
\end{array}
\qquad
 \nu_p(q_m)=\sum_{r\geq1}\left\lfloor\frac{m^2}{p^r}\right\rfloor
                         -m\nu_p(m).
\]
Thus $30\mid\sigma_m(m)q_m$ for $m\geq3$.
For $m=2$ the normalized coefficient is $(1+2^2)q_2=5\cdot6=30$,
and for $m=1$ it is one. Lemma~\ref{lem:transfer} with $M=30$ and $t=1$
proves the corollary.
\end{proof}

Both moduli in the corollaries are also optimal as uniform moduli:
$v_3-1=3$ and $z_4-1=30$. These values follow by multiplying
$\ex(x)$ by $1+x^3/2+O(x^6)$ and $1+5x^4/4+O(x^8)$, respectively.

\begin{corollary}\label{cor:sharpweights}
For any integers $t,w_2,w_3,\ldots$, let
$u_n=n![x^n]\exp(tx+\sum_{m\geq2}w_m x^{m^2}/m^m)$.
Then $u_n\in\mathbb Z$ and $u_n\equiv t^n\pmod6$ for all $n\geq0$.
Moreover, $u_n\equiv t^n\pmod{30}$ for every $n\geq0$
if and only if $5\mid w_2$.
\end{corollary}
\begin{proof}
Lemma~\ref{lem:six} gives integrality and the mod-six assertion.
The proof of Corollary~\ref{cor:weighted} gives $30\mid(m^2)!/m^m$
for $m\geq3$, so the sole condition in Lemma~\ref{lem:transfer}
for modulus thirty is $30\mid6w_2$, equivalently $5\mid w_2$.
The necessity is also visible in $u_4=t^4+6w_2$.
\end{proof}

For completeness, the source's values at $x=1$ in A226838 and A226839
also follow from the definitions. Both inner series have nonnegative
coefficients and converge at one: $m^{-m}\leq2^{-m}$ for $m\geq2$,
and $\sum_{m\geq1}1/m!=e-1$. Monotone convergence, applied to the
nonnegative coefficient expansion of the exponential, gives
\[
 \sum_{n\geq0}\frac{a_n}{n!}
   =\exp\!\left(\sum_{m\geq1}\frac1{m^m}\right),\qquad
 \sum_{n\geq0}\frac{v_n}{n!}=\exp(e-1).
\]

\section{Divisor-sum companions}

Write $s_n=\sigma_n(n)$ and $\sigma(n)=\sum_{d\mid n}d$ for $n\geq1$.
The linked entry A023887 records an identity conjectured by Velin Yanev
on 24 August 2017~\cite{divisors}. Its case $r=2$ is Fried's
Theorem~3~\cite{fried}; the same prime-power calculation proves all $r\geq1$.

\begin{corollary}\label{cor:radical}
For all integers $n,r\geq1$, with
$\operatorname{rad}(n)=\prod_{p\mid n}p$ and $\operatorname{rad}(1)=1$,
\[
 \sigma_r(n)=
 \frac{\sigma(n^r\operatorname{rad}(n)^{r-1})}
      {\sigma(\operatorname{rad}(n)^{r-1})}.
\]
\end{corollary}
\begin{proof}
For $n=\prod_{p\mid n}p^{a_p}$, independent choices of prime-power
divisors give
\[
 \sigma_r(n)=\prod_{p\mid n}\sum_{j=0}^{a_p}p^{rj}.
\]
The factor at $p$ on the right-hand side is
\[
 \frac{1+p+\cdots+p^{r(a_p+1)-1}}
      {1+p+\cdots+p^{r-1}}
   =\frac{p^{r(a_p+1)}-1}{p^r-1}
   =\sum_{j=0}^{a_p}p^{rj}.
\]
All denominators are positive. Multiplication proves the identity;
for $n=1$ both products are empty and equal one.
\end{proof}

\begin{corollary}\label{cor:divisors}
The following statements recorded in A023887 hold:
\begin{enumerate}
\item $s_n=\prod_{p^a\parallel n}(p^{n(a+1)}-1)/(p^n-1)$;
$s_n$ is odd exactly when $n$ is a square or twice a square.
\item As formal power series,
\[
 \sum_{n\geq1}s_nx^n=\sum_{d\geq1}\frac{(dx)^d}{1-(dx)^d}.
\]
For $D(x)=\exp(\sum_{n\geq1}s_nx^n/n)$, we have
\[
 D(x)=\prod_{d\geq1}(1-d^dx^d)^{-1/d},\qquad
 xD'(x)/D(x)=\sum_{n\geq1}s_nx^n.
\]
\item $1\leq s_n/n^n\leq1+n2^{-n}$, so $s_n\sim n^n$
and $\sum_{n\geq1}1/s_n$ converges.
\end{enumerate}
\end{corollary}
\begin{proof}
The product formula is the divisor factorization above with $r=n$.
For odd $n$, every divisor power is odd, so $s_n$ is odd exactly when
the number of divisors is odd, namely when every prime exponent is even.
For even $n$, the even divisors contribute zero modulo two, so the same
test applies to the odd part of $n$: it must be a square. According to
the parity of the exponent of two, this says square or twice a square.

Expanding the geometric series in (2), the coefficient of $x^n$ is
$\sum_{d\mid n}d^n=s_n$. Taking the formal logarithm of the product for
$D$ gives $\sum_{d,k\geq1}d^{dk}x^{dk}/(dk)$, whose coefficient at degree
$n$ is $s_n/n$. Formal differentiation proves the last equality.
This is the generating function of A023881~\cite{expanding}; the logarithmic
derivative statement uses $xD'/D$, not $D'/D$ without the factor $x$.

For (3), the divisor $n$ contributes $n^n$, and every proper divisor
is at most $n/2$. There are at most $n$ such divisors. Dividing by $n^n$
gives the displayed bound. The reciprocal series is bounded above by
$\sum_{n\geq1}n^{-n}<\infty$.
\end{proof}

The partition comparison in A226890 also follows by formal logarithms:
\[
 \prod_{d\geq1}(1-x^d)^{-1}
 =\exp\!\left(\sum_{n\geq1}\frac{\sigma(n)}{n}x^n\right),
\]
since the coefficient at $x^n$ in the logarithm is
$\sum_{d\mid n}1/(n/d)=\sigma(n)/n$.
Every formal sum and product here is finite in each coefficient.

The primality conjecture that $s_2=5$ is the only prime among the $s_n$
is a separate assertion in A023887 (Rayhan Ahmed, 19 June 2026).
It is not settled by these congruences or divisor identities. Ahmed's
question and Xiao's answer discuss the remaining rapidly growing
cyclotomic cases~\cite{primality}.

\section{Coverage and verification}

\begin{center}
\small
\begin{tabular}{@{}>{\raggedright\arraybackslash}p{.32\linewidth}>{\raggedright\arraybackslash}p{.59\linewidth}@{}}
\toprule
Source statement&Mathematical coverage\\\midrule
A226838, rev.\ 14, FORMULA&Proved for every $n\geq0$ in Theorem~\ref{thm:square}.\\
A226839, rev.\ 11, FORMULA&Proved for every $n\geq0$ in Corollary~\ref{cor:triangular}.\\
A226890, rev.\ 5, FORMULA&Proved for every $n\geq0$ in Corollary~\ref{cor:weighted}.\\
A226838 and A226839, $x=1$&Proved by monotone convergence.\\
A226890, partition comparison&Proved by formal logarithms.\\
A023887, rev.\ 53, Yanev FORMULA&Proved, Corollary~\ref{cor:radical}; $r=2$ reproved, prior proof: Fried (2025), Theorem~3.\\
A023887, established formulas&Proved, Corollary~\ref{cor:divisors}: product, parity, generating function, asymptotic and $xD'/D$ identity.\\
A023887, reciprocal sum&Convergence proved; decimal expansion not covered.\\
A023887, primality conjecture&Not covered.\\
Weighted mod-thirty criterion&Not posed; proved, Corollary~\ref{cor:sharpweights}.\\
\bottomrule
\end{tabular}
\end{center}

The Lean declaration \texttt{HannaA226838Congruence.result}
in \href{https://github.com/the-omega-institute/trureturing/blob/6974dba01daa1a25cc74b87a86c191e95fe2c681/D5/S1/Recurrence/Residue/HannaA226838Congruence.lean}{the square-exponent source file}
states, for the literal rational series \eqref{eq:square},
$\forall n\in\mathbb N\ \exists k\in\mathbb Z,\ n![x^n]A=1+6k$.
Its counterpart \texttt{HannaA226839Congruence.result} in
\href{https://github.com/the-omega-institute/trureturing/blob/6974dba01daa1a25cc74b87a86c191e95fe2c681/D5/S1/Recurrence/Residue/HannaA226839Congruence.lean}{the triangular-exponent source file}
states $\forall n\in\mathbb N\ \exists k\in\mathbb Z,\ n![x^n]B=1+3k$.
Both declarations include integrality and the index zero.
The declaration \texttt{YanevSigmaRadicalIdentity.result} in
\href{https://github.com/the-omega-institute/trureturing/blob/6974dba01daa1a25cc74b87a86c191e95fe2c681/D5/S3/Arith/YanevSigmaRadicalIdentity.lean}{the divisor-sum source file}
proves the multiplied-out natural-number equality in
Corollary~\ref{cor:radical} for positive $n,r$; positivity of the denominator
gives the quotient form stated here.
The transfer criterion, weighted extensions, modulo-thirty congruence,
optimality witnesses and remaining divisor identities have written proofs here.
The accompanying \texttt{verify.py} uses exact integer and rational
arithmetic to check finite coefficient ranges; its output specifies those ranges.
The accompanying \texttt{sibling-checks.py} checks bounded divisor identities
and weighted coefficient examples with exact arithmetic.

\begin{thebibliography}{8}\small
\bibitem{square}
Paul D. Hanna, \emph{E.g.f.: $\exp(\sum_{m\geq1}x^{m^2}/m^m)$},
OEIS A226838, revision 14 (20 June 2013), FORMULA and COMMENTS.
\url{https://oeis.org/A226838};
\href{https://oeis.org/history?seq=A226838}{revision history}.
\bibitem{triangle}
Paul D. Hanna, \emph{E.g.f.: $\exp(\sum_{m\geq1}x^{m(m+1)/2}/m!)$},
OEIS A226839, revision 11 (19 June 2013), FORMULA and COMMENTS.
\url{https://oeis.org/A226839}.
\bibitem{weighted}
Paul D. Hanna, \emph{E.g.f.: $\exp(\sum_{m\geq1}\sigma_m(m)x^{m^2}/m^m)$},
OEIS A226890, revision 5 (20 June 2013), FORMULA.
\url{https://oeis.org/A226890}.
\bibitem{divisors}
Olivier G\'erard and contributors, \emph{Sum of $n$th powers of divisors of $n$},
OEIS A023887, revision 53 (24 June 2026), COMMENTS and FORMULA.
The radical conjecture is attributed to Velin Yanev (24 August 2017).
\url{https://oeis.org/A023887}.
\bibitem{fried}
Sela Fried, \emph{Proofs of some conjectures of Yanev} (2025),
Theorem~3, p.~3; linked from OEIS A001157, FORMULA, 1 December 2025.
\url{https://oeis.org/A001157/a001157.pdf}.
\bibitem{expanding}
Olivier G\'erard and contributors, \emph{Number of partitions in expanding space},
OEIS A023881, revision 35 (8 September 2022), FORMULA;
exponential formula by Michael Somos (2006), product by Paul D. Hanna (2011).
\url{https://oeis.org/A023881}.
\bibitem{primality}
Rayhan Ahmed, \emph{Is $\sigma_2(2)=5$ the only prime value of the diagonal
divisor power sum $\sigma_n(n)$?}, MathOverflow question 512305,
14 June 2026; answer by Stanley Yao Xiao, 15 June 2026.
\url{https://mathoverflow.net/q/512305}; \url{https://mathoverflow.net/a/512320}.
\bibitem{carlitz}
L. Carlitz, \emph{Some congruences for the Bell polynomials},
Pacific Journal of Mathematics \textbf{11}(4) (1961), 1215--1222;
\S2, (2.1)--(2.2), p.~1216. \url{https://doi.org/10.2140/pjm.1961.11.1215}.
\end{thebibliography}

\noindent\small Licensed under Creative Commons Attribution 4.0 International
(CC BY 4.0).
\end{document}
